% chktex-file 9
% chktex-file 17

本研究以蒙地卡羅模擬計算理想條件下（不含雜散光與量測雜訊）的角度分布，作為實驗對照。每顆光子自比色槽入口射入，先依
\[
s=-\frac{\ln\xi}{\mu_s}
\]
前進隨機距離，其中 $\xi$ 為 $0$ 至 $1$ 之間的亂數，$\mu_s$ 為散射係數；碰到微球後，依含粒徑分布的 Mie 相位函數決定新方向，粒徑分布的變異係數分別採用 $0.492\,\mu\mathrm{m}$ 的 $2.1\%$ 與 $2.16\,\mu\mathrm{m}$ 的 $3.1\%$。光子碰到比色槽壁時，依 Fresnel 公式處理水、PMMA（折射率 $1.489$）與空氣界面的反射及透射，包含槽壁內多次反射與全反射；穿出後再依 Snell 定律折射，記錄其落入旋轉臂的 $1^\circ$ 角度格。

每個 $\tau$ 模擬 $5{,}000$ 萬顆光子；光子偏離光束平面超過 $150\,\mathrm{mm}$ 後停止追蹤。模擬假設雷射為非偏振光，旋轉臂半徑取實測值 $110\,\mathrm{mm}$，並以遮光管收光半角 $15.3^\circ$（管長 $11\,\mathrm{mm}$、開口寬 $6\,\mathrm{mm}$）限制感測器收光。每個條件使用 $6$ 組不同亂數種子，取模擬平均值，再以和實驗相同的方式計算 $C$。驗證時，移除比色槽、令 $\tau\to0$ 並只計入單次散射光子，模擬結果與 Mie 理論一致；另以獨立重寫的程式計算六個條件，結果在統計誤差內一致。

模擬方式：
\subsection*{(一) 前置準備}
\begin{itemize}
    \item[1.] 計算出粒子的 $S_1(\theta),\ S_2(\theta)$ 存成查表
    \item[2.] 計算出粒子的散射、吸收截面$\left(\sigma_s,\ \sigma_a\right)$
    \item[3.] 計算出粒子的散射、吸收係數$\left(\mu_s,\ \mu_a\right)$
    \item[4.] 將穆勒矩陣的 $S_{11}$ 轉成累積分布函數 $CDF$ 
\end{itemize}
\subsection*{(二) 模擬內容}
\subsection*{1. 抽取自由路徑}
按照自由路徑分布公式$s = -\lambda\ln\left(\zeta\right),\ \zeta \in (0, 1]$